\section{Diagram 6}\label{sec:d6}

\subsection{The diagram and the butterfly chart}

Diagram~6 is the \emph{butterfly} of \cite[\S7.4]{AK}, a $z$-triangle and a $w$-triangle (the
\emph{wings}) with one common vertex, with both wings circled.  In the numbering of
\cite[\S8.6]{AK} the $z$-wing $\{1,2,5\}$ has the $z$-strokes $12,15,25$ and $z$-circles at $1,2$;
the $w$-wing $\{3,4,5\}$ has the $w$-strokes $34,35,45$ and $w$-circles at $3,4$; vertex $5$ is
not circled, and no stroke joins $\{1,2\}$ to $\{3,4\}$.
The butterfly has three circlings, Diagrams~6--8, and $5$ is never circled
\cite[\S7.4]{AK} (Fact~\ref{fact:fig11}); in Diagram~6 both wings are circled, and a circled
$z$-wing has $z$-circles exactly at $1,2$ (Rules~1a, 1d), and likewise for $w$.

Permutations of the bodies (with the masses) and $z\leftrightarrow w$ are symmetries of
\eqref{eq:system}.  The thirty coloured labels of Diagram~6 (a common vertex and an ordered pair
of wings) are images of this numbering, and $z\leftrightarrow w$ composed with $(13)(24)$
preserves it.  The chart below uses rotation invariants only, so the gauge $z_{12}=w_{12}$ of
$A_m$ plays no role, and it suffices to treat this numbering at arbitrary positive masses.

\begin{definition}\label{d6-def-chart}
The \emph{butterfly chart} has the $21$ coordinates $t$, $Z_1,\dots,Z_5$, $W_1,\dots,W_5$,
$D_{kl}$, related to \eqref{eq:system} by $(z_k,w_k)=(t^{-1}Z_k,tW_k)$ for $k=1,2$,
$(z_k,w_k)=(tZ_k,t^{-1}W_k)$ for $k=3,4$, $(z_5,w_5)=(tZ_5,tW_5)$, $\delta_{kl}=tD_{kl}$ on the
\emph{crossed pairs} $13,14,23,24$ and $\delta_{kl}=D_{kl}$ on the \emph{wing pairs}
$12,15,25,34,35,45$.  Its \emph{rows} are the $z$- and $w$-equations
$z_k-\sum_{l\ne k}m_l\delta_{kl}^3(z_k-z_l)=0$ (likewise for $w$) of body $k$, multiplied by
$(t,t^{-1})$ for $k=1,2$, by $(t^{-1},t)$ for $k=3,4$ and by $(t,t)$ for $k=5$ and denoted
$R^z_k,R^w_k$; the ten distance equations $\delta_{kl}^2z_{kl}w_{kl}-1=0$; and the gauges
$Z_2-Z_1-1=0$, $W_4-W_3-1=0$.  We write $F=0$ for these $22$ rows and
$\mathcal B=F^{-1}(0)\cap\{t=0\}$ for the \emph{boundary fibre}.  ($Z_k,W_k$ are not the terms
$Z_{kl}=z_{kl}\delta_{kl}^3$ of \cite[\S3]{AK}.)
\end{definition}

\begin{lemma}[computer-assisted]\label{d6-lem-chart}
Let $m\in\Rp^5$.
\begin{romanlist}
  \item The rows are polynomials in the $21$ coordinates, of degree at most $5$ in $t$.
  \item Let $(z,w,\delta)\in A_m$ (so $z_{12}w_{34}\ne0$), choose $t$ with
  $t^2=(z_{12}w_{34})^{-1}$ and put $\nu=(tz_{12})^{-1}$.  The chart point of the rotated
  solution $(\nu z,\nu^{-1}w,\delta)$ has $Z_k=z_k/z_{12}$, $W_k=w_kz_{12}$ $(k=1,2)$,
  $Z_k=z_kw_{34}$, $W_k=w_k/w_{34}$ $(k=3,4)$, $Z_5=z_5w_{34}$, $W_5=w_5z_{12}$,
  $D_{kl}=\delta_{kl}/t$ on crossed pairs and $D_{kl}=\delta_{kl}$ on wing pairs, and satisfies
  $F=0$: each coordinate row is $t^{\pm1}\nu^{\pm1}$ times an equation of \eqref{eq:system}, the
  distance rows are those of \eqref{eq:system}, and both gauges hold.
  \item At $t=0$ the crossed distance rows read $D_{kl}^2Z_kW_l=-1$, and the wing distance rows
  read $D_{12}^2(Z_2-Z_1)(W_2-W_1)=1$, $-D_{15}^2Z_1(W_5-W_1)=1$, $-D_{25}^2Z_2(W_5-W_2)=1$,
  $D_{34}^2(Z_4-Z_3)(W_4-W_3)=1$, $-D_{35}^2(Z_5-Z_3)W_3=1$, $-D_{45}^2(Z_5-Z_4)W_4=1$.  So
  every $D_{kl}$ is nonzero on $\mathcal B$.
  \item Identically,
  $m_1R^z_1+m_2R^z_2+m_5R^z_5+t^2(m_3R^z_3+m_4R^z_4)=m_1Z_1+m_2Z_2+t^2(m_3Z_3+m_4Z_4+m_5Z_5)$,
  and likewise after $z\leftrightarrow w$, $1\leftrightarrow3$, $2\leftrightarrow4$.  Hence on
  $\mathcal B$: $Z_1=-m_2/(m_1+m_2)$, $Z_2=m_1/(m_1+m_2)$, $W_3=-m_4/(m_3+m_4)$,
  $W_4=m_3/(m_3+m_4)$.
  \item At $t=0$ the rows $R^z_1,R^z_2,R^z_5,R^w_1,R^w_2$, the distance rows of $12,15,25$ and
  the first gauge involve only the \emph{$z$-wing coordinates}
  $Z_1,Z_2,W_1,W_2,W_5,D_{12},D_{15},D_{25}$; the remaining rows except the crossed distance rows
  involve only the \emph{$w$-wing coordinates} $Z_3,Z_4,Z_5,W_3,W_4,D_{34},D_{35},D_{45}$; and
  $D_{13},D_{14},D_{23},D_{24}$ occur only in the crossed distance rows.
\end{romanlist}
\end{lemma}
We check (i)--(iv) by exact expansion over $\mathbb Q$, with the masses and the coordinates as
indeterminates; for (ii) we substitute the chart point of the rotated solution into the $22$
rows (Appendix~\ref{app:verification}).  Item (v) holds because every term coupling the two
wings carries a positive power of $t$.

\subsection{Coverage}

\begin{proposition}[Coverage]\label{prop:d6-cover}
Let $m\in\Rp^5$ and let $\gamma(s)$, $s\to0$, be a Diagram-6 germ, in the numbering above, on a
curve in $A_m$.  Along $\gamma$ define $t$ and $\nu$ as in Lemma~\ref{d6-lem-chart}(ii), and
let $\tilde\gamma(s)$ be the chart point of the rotated solution.  Then $t\to0$, every chart
coordinate is bounded along $\tilde\gamma$, and $\tilde\gamma(s)$ converges to a point
$p\in\mathcal B$ at which $Z_1,Z_2,W_3,W_4$ and all ten $D_{kl}$ are nonzero.
\end{proposition}

\begin{proof}
Normalize as in \cite[p.~547]{AK} so that the largest components of $\mathcal Z$ and $\mathcal W$
both have modulus $\epsilon^{-2}$; by Lemma~\ref{lem:germ} this is realized along $\gamma$, with
$\epsilon$ comparable to a positive power of $s$.  Let $\ord$ be the order in $\epsilon$
($\ord f>\alpha$ means $f\epsilon^{-\alpha}\to0$).  The chart coordinates are invariant under
$z\mapsto az$, $w\mapsto a^{-1}w$, so we compute in this frame.

By the circles, $\ord z_1=\ord z_2=\ord w_3=\ord w_4=-2$, and the uncircled
$z_3,z_4,z_5,w_1,w_2,w_5$ have order $>-2$; so $\ord z_{15}=\ord z_{25}=-2$.  A $z$-stroke $kl$
means $\ord(z_{kl}\delta_{kl}^3)=\ord(z_{kl}^{-1/2}w_{kl}^{-3/2})=-2$, that is,
$\ord z_{kl}+3\ord w_{kl}=4$.  So $\ord w_{15}=\ord w_{25}=2$, hence $\ord w_{12}\ge2$, and the
stroke $12$ gives $\ord z_{12}\le-2$; with $\ord z_{12}\ge-2$ (Estimate~1,
Fact~\ref{fact:rules}) we get $\ord z_{12}=-2$, $\ord w_{12}=2$.  Thus $r_{kl}^2=z_{kl}w_{kl}$ and
$\delta_{kl}$ have order $0$ on the three $z$-wing pairs (maximal edges, $r\approx1$ as in
Estimate~2), and symmetrically on the $w$-wing.  On a crossed pair $z_{kl}$ and $w_{kl}$ are
differences of a circled and an uncircled coordinate, so both have order $-2$ and
$\ord\delta_{kl}=2$.  In the $w$-equation $w_1=\sum_{l\ne1}m_l\delta_{1l}^3(w_1-w_l)$ the terms
have orders $2$ ($l=2,5$) and $4$ ($l=3,4$); so $\ord w_1\ge2$, likewise
$\ord w_2,\ord z_3,\ord z_4\ge2$, and $w_5=w_1+w_{15}$, $z_5=z_3+z_{35}$ have order $\ge2$.
(These are the estimates $A_w,B_z,m_1w_1+m_2w_2,w_5,B_w,A_z\preccurlyeq\epsilon^2$ of
\cite[\S8.6]{AK}.)

Finally $\ord t=-\frac12(\ord z_{12}+\ord w_{34})=2$ and $\ord\nu=-\ord t-\ord z_{12}=0$.  By
Lemma~\ref{d6-lem-chart}(ii), $Z_1,Z_2,W_3,W_4$ and all ten $D_{kl}$ have order exactly $0$
(crossed pairs: $2+\frac12(-2-2)=0$), and $W_1,W_2,W_5,Z_3,Z_4,Z_5$ have order $\ge0$.  Bounded
Puiseux series in $s$ converge, so $\tilde\gamma(s)\to p\in\mathcal B$, and the coordinates of
order exactly $0$ have nonzero limits.
\end{proof}
In particular $D_{12}=\delta_{12}$ and $D_{34}=\delta_{34}$ have finite nonzero limits along every
Diagram-6 germ, at every mass.

\subsection{The first condition}

Put $\Sigma=\sum_{k\in\{1,2\},\,l\in\{3,4\}}m_km_lD_{kl}^3Z_kW_l$.

\begin{lemma}[computer-assisted]\label{d6-lem-torque}
There is a polynomial $Q$ in the chart coordinates, with coefficients polynomial in $m$, such
that identically
\begin{equation}\label{d6-eq-torque}
  \sum_{k=1,2}m_k\bigl(W_kR^z_k-Z_kR^w_k\bigr)+m_5W_5R^z_5=-t\,\Sigma+t^2Q.
\end{equation}
\end{lemma}
Expanding the left-hand side of \eqref{d6-eq-torque} in exact arithmetic over $\mathbb Q$, with
the masses and the chart coordinates as indeterminates, we find that its coefficients of $t^0$
and $t^1$ are $0$ and $-\Sigma$ (Appendix~\ref{app:verification}).
With $\rho^z_k,\rho^w_k$ the left-hand sides of \eqref{eq:system}, the pair torque
$\sum_{k=1,2}m_k(w_k\rho^z_k-z_k\rho^w_k)$ equals
$\sum_{k=1,2}\sum_{l=3,4,5}m_km_l\delta_{kl}^3(w_kz_l-z_kw_l)$ (the terms with $l\in\{1,2\}$
cancel), which is $2i$ times (31) of \cite[\S8.6]{AK}.  In the chart it is the first sum in
\eqref{d6-eq-torque}; its terms with $l=3,4$ are $-t\Sigma+t^5\sum m_km_lD_{kl}^3W_kZ_l$, those with
$l=5$ are $\sum_{k=1,2}m_km_5D_{k5}^3(t^2W_kZ_5-Z_kW_5)$, and
$R^z_5-(m_1D_{15}^3Z_1+m_2D_{25}^3Z_2)$ is divisible by $t^2$.

\begin{proposition}[First condition]\label{prop:d6-first}
Let $m\in\Rp^5$, let $\gamma$ be a Diagram-6 germ in the numbering above, and let $p$ be its
limit (Proposition~\ref{prop:d6-cover}).  Then:
\begin{romanlist}
  \item $\Sigma(p)=0$, and $\Sigma(p)=-(m_1m_3D_{13}+m_1m_4D_{14}+m_2m_3D_{23}+m_2m_4D_{24})(p)$;
  \item $m_1=m_2$ and $m_3=m_4$; this is relation (32) of \cite[\S8.6]{AK};
  \item $(D_{13},D_{14},D_{23},D_{24})(p)=(2i\eta_{13},\,2\eta_{14},\,2\eta_{23},\,2i\eta_{24})$
  with signs $\eta_{kl}\in\{\pm1\}$ satisfying $\eta_{14}=-\eta_{23}$ and $\eta_{13}=-\eta_{24}$.
\end{romanlist}
In another numbering, wings $\{a,b,c\}$ and $\{d,e,c\}$ force $m_a=m_b$ and $m_d=m_e$.
\end{proposition}

\begin{proof}
All rows vanish along $\tilde\gamma$, so \eqref{d6-eq-torque} gives $t\Sigma=t^2Q$, and $t\ne0$
gives $\Sigma=tQ$.  By Proposition~\ref{prop:d6-cover}, $Q$ is bounded and $t\to0$, so
$\Sigma(p)=0$.  By Lemma~\ref{d6-lem-chart}(iii), $D_{kl}^2Z_kW_l=-1$ at $p$ for the crossed
pairs, so $\Sigma(p)=\sum m_km_lD_{kl}(D_{kl}^2Z_kW_l)=-\sum m_km_lD_{kl}$; this is (i).

By Lemma~\ref{d6-lem-chart}(iv), $Z_1W_3=m_2m_4/s_0^2$, $Z_1W_4=-m_2m_3/s_0^2$,
$Z_2W_3=-m_1m_4/s_0^2$ and $Z_2W_4=m_1m_3/s_0^2$ at $p$, where $s_0>0$,
$s_0^2=(m_1+m_2)(m_3+m_4)$.  Write $m_k=\mu_k^2$ with $\mu_k>0$.  Then $D_{kl}^2=-1/(Z_kW_l)$ gives
\[
  D_{13}=\frac{i\eta_{13}s_0}{\mu_2\mu_4},\quad D_{14}=\frac{\eta_{14}s_0}{\mu_2\mu_3},\quad
  D_{23}=\frac{\eta_{23}s_0}{\mu_1\mu_4},\quad D_{24}=\frac{i\eta_{24}s_0}{\mu_1\mu_3}
\]
with signs $\eta_{kl}\in\{\pm1\}$, and
\[
  \Sigma(p)=-\frac{s_0}{\mu_1\mu_2\mu_3\mu_4}\bigl(P_{\mathrm{re}}+iP_{\mathrm{im}}\bigr),\qquad
  P_{\mathrm{re}}=\eta_{14}\mu_1^3\mu_4^3+\eta_{23}\mu_2^3\mu_3^3,\quad
  P_{\mathrm{im}}=\eta_{13}\mu_1^3\mu_3^3+\eta_{24}\mu_2^3\mu_4^3.
\]
Up to the factor $\pm i$, $P_{\mathrm{re}}+iP_{\mathrm{im}}$ is the left-hand side of (12) of
\cite[\S5.1]{AK} written in the $\mu_k$, and each of its eight sign choices arises from exactly two
of the sixteen sign classes $\eta$.
Since the $\mu_k$ and $s_0$ are real, $\Sigma(p)=0$ forces $P_{\mathrm{re}}=P_{\mathrm{im}}=0$.  If
$\eta_{14}=\eta_{23}$, then $P_{\mathrm{re}}\ne0$, and if $\eta_{13}=\eta_{24}$, then
$P_{\mathrm{im}}\ne0$; this excludes twelve classes.  In the other four,
$P_{\mathrm{re}}=\eta_{14}(x-y)(x^2+xy+y^2)$ with $x=\mu_1\mu_4$, $y=\mu_2\mu_3$ and
$x^2+xy+y^2>0$, so $\mu_1\mu_4=\mu_2\mu_3$; likewise $P_{\mathrm{im}}=0$ gives
$\mu_1\mu_3=\mu_2\mu_4$.  Then
$\mu_3(\mu_1\mu_4-\mu_2\mu_3)-\mu_4(\mu_1\mu_3-\mu_2\mu_4)=\mu_2(\mu_4-\mu_3)(\mu_4+\mu_3)$
vanishes, so $\mu_3=\mu_4$ and then $\mu_1=\mu_2$: this is (ii).  With $m_1=m_2$ and $m_3=m_4$
the displayed values of $D_{kl}$ become those of (iii).
\end{proof}
At $m_1=m_2$, $m_3=m_4$ the expression $\Sigma(p)$ vanishes for exactly the four classes of (iii).
Albouy and Kaloshin \cite[\S8.6]{AK} obtain (32) ``as in 5.1'' from the estimate
$S\preccurlyeq\epsilon^4$ in (31); the proof above makes this exact in the chart.

\subsection{The finite wing atlas}

From now on $m=(a,a,b,b,c)$ with $a,b,c>0$.

\begin{lemma}[computer-assisted]\label{d6-lem-atlas}
The points of $\mathcal B$ are exactly the points with $Z_1=W_3=-\frac12$, $Z_2=W_4=\frac12$,
\begin{gather*}
  (D_{12},D_{15},D_{25})=(d,\lambda d,\omega\lambda d),\\
  W_1=-d(a+2c\lambda),\quad
  W_2=d(a+2c\omega\lambda),\quad W_5=ad\bigl(4\lambda^{-2}-1\bigr),\\
  (D_{34},D_{35},D_{45})=(e,\kappa e,\omega'\kappa e),\\
  Z_3=-e(b+2c\kappa),\quad
  Z_4=e(b+2c\omega'\kappa),\quad Z_5=be\bigl(4\kappa^{-2}-1\bigr),\\
  (D_{13},D_{14},D_{23},D_{24})=(2i\eta_{13},\,2\eta_{14},\,2\eta_{23},\,2i\eta_{24}),\quad
  \eta_{kl}\in\{\pm1\},
\end{gather*}
where
\begin{equation}\label{d6-eq-wing}
  \omega^3=1,\quad\lambda^2=2(1+\omega),\quad(2a+c\lambda^3)d^3=1,\qquad
  \omega'^3=1,\quad\kappa^2=2(1+\omega'),\quad(2b+c\kappa^3)e^3=1.
\end{equation}
For each $\omega$ both roots $\lambda$ are nonzero.  The scale denominator $2a+c\lambda^3$
vanishes only for $(\omega,\lambda)=(1,-2)$ and $a=4c$: for $\omega=1$, $\lambda=\pm2$; for
$\omega\ne1$, $\lambda^6=-8$ and the two denominators have product $4a^2+8c^2$.  Hence a wing
has $18$ classes $(\omega,\lambda,d)$ if $a\ne4c$ and $15$ if $a=4c$, and $\mathcal B$ has
$18\cdot18\cdot16=5184$ points if $a\ne4c\ne b$, $4320$ if exactly one of $a,b$ equals $4c$, and
$3600$ if $a=b=4c$.
\end{lemma}
The proof is an exact computation over the field $\mathbb Q(i)(a,b,c)$, with the rows reduced
modulo \eqref{d6-eq-wing} by a Gr\"obner basis, and follows the argument below
(Appendix~\ref{app:verification}).  By Lemma~\ref{d6-lem-chart}(iv),(v), $Z_1=-\frac12$,
$Z_2=\frac12$, and $\mathcal B$ is the product of the two wing fibres and the crossed roots.
On the $z$-wing put $\lambda=D_{15}/D_{12}$ and $\omega=D_{25}/D_{15}$.  At $t=0$,
$2(R^z_1+R^z_2)/c=D_{15}^3-D_{25}^3$ gives $\omega^3=1$; the wing distance rows give
$D_{12}^{-2}=W_2-W_1=2D_{15}^{-2}+2D_{25}^{-2}$,
that is, $\lambda^2=2(1+\omega)$; and $2R^z_1=(2a+c\lambda^3)d^3-1$.  The row $R^w_1$ and the
distance rows of $12,15$ are linear in $(W_1,W_2,W_5)$ with determinant $d^4\lambda^2/2\ne0$;
conversely the displayed point satisfies all eight $z$-wing rows modulo \eqref{d6-eq-wing}.  The
$w$-wing follows by the wing exchange below.

Germs at infinity are complex (Remark~\ref{rem:levels}), so all eighteen classes are needed, not
only the positive real class $(\omega,\lambda)=(1,2)$, $D_{15}=D_{25}=2D_{12}$.

\begin{definition}\label{d6-def-labels}
A \emph{label} is a point of $\mathcal B$, that is, a choice of $(\omega,\lambda,d)$,
$(\omega',\kappa,e)$ and $\eta$ in Lemma~\ref{d6-lem-atlas}.  The wing \emph{shapes}
$(\omega,\lambda)$ and $(\omega',\kappa)$ are \emph{Gaussian} if $\omega=1$ (so $\lambda=\pm2$) and
\emph{primitive} otherwise (so $\lambda^4-2\lambda^2+4=0$); each wing has two Gaussian and four
primitive shapes.  A label is \emph{corank-one} if $\eta_{14}=-\eta_{23}$ and
$\eta_{13}=-\eta_{24}$, and \emph{regular} otherwise.  For $a\ne4c\ne b$ there are
$12\cdot324=3888$ regular and $4\cdot324=1296$ corank-one labels.
\end{definition}

\emph{Symmetries.}  At the masses $(a,a,b,b,c)$, $F^{-1}(0)$ is preserved by the exchange of $1$
and $2$ combined with $t\mapsto it$, $(Z_1,Z_2,W_1,W_2,W_5)\mapsto(-Z_2,-Z_1,-W_2,-W_1,-W_5)$ and
$D_{kl}\mapsto-iD_{kl}$ on crossed pairs, by $(t,D_{13},D_{14},D_{23},D_{24})\mapsto
-(t,D_{13},D_{14},D_{23},D_{24})$, and by complex conjugation; the \emph{wing exchange}
($z\leftrightarrow w$, $1\leftrightarrow3$, $2\leftrightarrow4$) maps the chart at $(a,a,b,b,c)$ to
that at $(b,b,a,a,c)$.  On the $576$ \emph{states} ($z$-wing shape, $w$-wing shape, $\eta$) these
maps have $29$ orbits, $21$ regular and $8$ corank-one (Appendix~\ref{app:verification}).  The
computations below express the masses through the scales, $a=\frac12(d^{-3}-c\lambda^3)$,
$b=\frac12(e^{-3}-c\kappa^3)$, so one symbolic computation covers the nine choices of $(d,e)$ of a
state at every positive $(a,b,c)$.

\begin{remark}[Regular labels]\label{d6-rem-regular}
Regular labels are already excluded by Proposition~\ref{prop:d6-first}(iii).  Independently, $DF$
has rank $21$ at every regular label, which is therefore an isolated point of $F^{-1}(0)$: at one
state of each regular orbit the minor omitting $R^z_1$ is a unit in $\{\pm1,\pm i,\pm1\pm i\}$
times $-\frac{81}{128}a^2b\,d^2e^4H(\lambda)H(\kappa)(2a+4c\lambda)/(2a+c\lambda^3)$, where
$H(x)=x^3(x^2-2)^2((x^2-2)^2+4)$, and this is nonzero on finite labels.  The minors are computed
in exact arithmetic over the field $\mathbb Q(i,\sqrt2,\sqrt3)(c,d,e)$, and their nonvanishing is
checked with resultants and norms (Appendix~\ref{app:verification}).
\end{remark}

\begin{lemma}[computer-assisted]\label{d6-lem-corank}
Let $p$ be a corank-one label.
\begin{romanlist}
  \item The $20$ rows of $F$ other than $R^w_2$ and $R^w_4$, together with the coordinate $t$,
  have invertible derivative at $p$.  Hence near $p$ the common zeros of these $20$ rows form a
  smooth curve $C_p$, the graph of an analytic map of $t$.
  \item Along $C_p$, $R^w_2=c_kt^k+O(t^{k+1})$ with $c_k\ne0$, where $k$ and $c_k$ are as in
  Table~\ref{d6-tab-corank}.
  \item Consequently $p$ is an isolated point of $F^{-1}(0)$.
\end{romanlist}
\end{lemma}

\begin{table}[ht]
\centering
\small
\begin{tabular}{p{3.3cm}p{0.9cm}p{0.3cm}p{6.6cm}}
\hline
wing shapes ($z$, $w$) & labels & $k$ & $c_k$ is a nonzero multiple of \\
\hline
Gaussian, Gaussian & $144$ & $4$ &
$Q_{2,2}=-\frac{2a+2b+c}{60(a+4c)(b+4c)}$, $Q_{2,-2}=\frac{2a+2b+c}{4(a+4c)(b-4c)}$,
$Q_{-2,-2}=\frac{2a+2b+c}{68(a-4c)(b-4c)}$ for $(\lambda,\kappa)=(2,2),(2,-2),(-2,-2)$;
$(-2,2)$ by wing exchange \\
one Gaussian, one primitive & $576$ & $3$ &
$P=A\kappa^3-4A\kappa+B\kappa^3-B\kappa^2-4B\kappa+B$ for $\lambda=\pm2$, $\kappa$ primitive;
product over the primitive $\kappa$: $9(8A^2+16AB+9B^2)^2$ \\
primitive, primitive: orbit of $(\lambda_0,\lambda_0)$ & $288$ & $2$ &
$1-2\kappa^3$; product over the primitive $\kappa$: $1089$ \\
primitive, primitive: orbits of $(\lambda_0,-\lambda_0)$ & $288$ & $2$ &
$2B/(2B+\kappa^3)$; product over the primitive $\kappa$: $B^4/(B^2+2)^2$ \\
\hline
\end{tabular}
\caption{Order $k$ and leading coefficient of $R^w_2$ along $C_p$ at the corank-one labels
(Lemma~\ref{d6-lem-corank}), at one representative state of each orbit.  Counts are for
$a\ne4c\ne b$; $A=a/c$, $B=b/c$; $\lambda_0=\frac12(\sqrt6+i\sqrt2)$ is the primitive ratio for
$\omega=\frac12(-1+i\sqrt3)$.  In the first row $R^w_4$ has the same order.}\label{d6-tab-corank}
\end{table}

The proof of (i) and (ii) is an exact computation at one representative state of each of the
eight corank-one orbits, over the field $\mathbb Q(i)(c,d,e)$ if both wings are Gaussian and over
its extension by a root $\kappa$ of $\kappa^4-2\kappa^2+4$ otherwise
(Appendix~\ref{app:verification}).  In (i) the derivative is reduced to the Schur complement of
the diagonal block of the distance rows.  For two Gaussian wings its determinant (rows multiplied
by $d^3e^3$) is $82944\,d^{61}e^{59}(cd^3-\frac{\lambda}{16})(ce^3-\frac{\kappa}{16})$, with
$cd^3-\frac{\lambda}{16}=-\frac{\lambda a}{8(2a+c\lambda^3)}\ne0$.  Otherwise the complement, after
removal of the unit row of $t$, is block diagonal with one $5\times5$ block per wing
(Lemma~\ref{d6-lem-chart}(v)); we check that an off-diagonal block vanishes and evaluate both
diagonal blocks.  For one Gaussian and one primitive wing the two determinants are nonzero
multiples of $8cd^3-\frac{\lambda}{2}=-\lambda A/(2A+\lambda^3)$ and of a polynomial in $ce^3$ and
$\kappa$ whose product over the primitive $\kappa$ is
$(\frac9{16})^4\cdot64B^4(B^4-8B^2+64)/(B^2+2)^4$, where $B^4-8B^2+64=(B^2-4)^2+48$.  This
product and those in Table~\ref{d6-tab-corank} are obtained from resultants with
$\kappa^4-2\kappa^2+4$.  For (ii) the jets of $C_p$ are computed by solving the $20$ rows of (i)
order by order in $t$, and are checked by substitution into all $22$ rows.  Item (iii) follows
from (i) and (ii): near $p$, $F^{-1}(0)\subset C_p$, and on $C_p$ the row $R^w_2$ vanishes only at
$t=0$, that is, at $p$.  With Remark~\ref{d6-rem-regular}, every label is thus an isolated point
of $F^{-1}(0)$; for $a\ne4c\ne b$ these are $3888+1296=5184$ labels.

\subsection{Exclusion}

\begin{theorem}\label{thm:d6}
Let $m\in\Rp^5$.  No curve in $A_m$ has a Diagram-6 germ, for any numbering of the bodies and
either colouring.  In particular the case $r=6$ of Theorem~\ref{thm:exclusion}(a) holds, at the
complex level.
\end{theorem}

\begin{proof}
By the symmetries described at the beginning of this section we may use the numbering above.  Let
$\gamma(s)$ be a Diagram-6 germ on a curve in $A_m$.  By Proposition~\ref{prop:d6-cover} the
chart germ $\tilde\gamma(s)$ has $t(s)\to0$ with $t(s)\ne0$ for $s\ne0$, bounded coordinates,
and a limit $p\in\mathcal B$.  By Proposition~\ref{prop:d6-first}, $m=(a,a,b,b,c)$ and the
crossed signs of $p$ satisfy $\eta_{14}=-\eta_{23}$, $\eta_{13}=-\eta_{24}$; by
Lemma~\ref{d6-lem-atlas}, $p$ is a label, hence a corank-one label.  By
Lemma~\ref{d6-lem-corank}, $p$ is an isolated point of $F^{-1}(0)$.  The coordinates of
$\tilde\gamma$ are Puiseux series in $s$ without negative exponents, so after a substitution
$s=\sigma^N$ they are convergent power series in $\sigma$, and $\sigma\mapsto\tilde\gamma(\sigma^N)$
is a holomorphic map of a disc into $F^{-1}(0)$ with value $p$ at $\sigma=0$.  Since $p$ is
isolated, this map is constant near $\sigma=0$, contradicting $t\not\equiv0$.
\end{proof}

The argument covers every ramification, every crossed sign class and the faces $a=4c$, $b=4c$,
since Lemma~\ref{d6-lem-atlas} describes the whole fibre $\mathcal B$.

\begin{remark}\label{d6-rem-projective}
On the face $a=4c$ the wing class $(\omega,\lambda)=(1,-2)$ has no finite scale; it lies at
$D_{12}=\infty$.  The proof of Theorem~\ref{thm:d6} never meets this class: by
Proposition~\ref{prop:d6-cover}, $D_{12}$ and $D_{34}$ have finite nonzero limits along every
Diagram-6 germ.  In a one-wing projective chart with $h=1/D_{12}$, and on $a=b=4c$ in a
double-projective ratio chart, one can show that no analytic branch through a point of
$\{t=h=0\}$ leaves the divisor $h=0$, though not that these points are isolated
(Appendix~\ref{app:verification}).  We do not use this.
\end{remark}
